\documentclass[11pt]{article}
\usepackage{amsmath,amssymb,amsthm}
\usepackage[margin=1in]{geometry}
\usepackage{hyperref}

\newtheorem{proposition}{Proposition}
\newtheorem{lemma}{Lemma}
\newtheorem{theorem}{Theorem}
\newtheorem{corollary}{Corollary}
\newtheorem{remark}{Remark}
\newcommand{\R}{\mathbb{R}}
\newcommand{\E}{\mathbb{E}}
\newcommand{\dd}{\,\mathrm{d}}
\newcommand{\Dv}{D\bm{v}}
\newcommand{\bx}{\bm{x}}
\newcommand{\by}{\bm{y}}
\newcommand{\bz}{\bm{z}}
\newcommand{\bv}{\bm{v}}
\newcommand{\bp}{\bm{p}}
\usepackage{bm}

\title{Corrected Derivation of the Quasi-Linearized Cole-Hopf\\
Transformation for Viscous Hamilton-Jacobi PDEs}
\author{Mathematical Supplement to ICML 2026 Submission}
\date{\today}

\begin{document}
\maketitle

\begin{abstract}
We provide a corrected and complete derivation of the Cole-Hopf-like
transformation proposed for linearizing the viscous Hamilton-Jacobi (HJ) PDE.
We show that the transformation is \emph{exact} only when the Hamiltonian
is quadratic in the co-state ($H = \tfrac{1}{2}|\bp|^2$), and derive the
precise residual that arises for general Hamiltonians.  We then formulate
the quasi-linearization iteration and establish the Gaussian expectation
formulas for the value function and its spatial gradient.
\end{abstract}

%% ==========================================================================
\section{Setup}
%% ==========================================================================

Consider the viscous HJ initial-value problem:
\begin{equation}\label{eq:hjvisc}
  \bv_t^\delta + H(t;\bx, D\bv^\delta) = \frac{\delta}{2}\,\Delta \bv^\delta
  \quad\text{in } \R^n\times(0,T],
  \qquad
  \bv^\delta(0;\bx) = g(\bx),
\end{equation}
where $\delta > 0$ is the viscosity parameter, $H:\R\times\R^n\times\R^n\to\R$
is the Hamiltonian, and $g:\R^n\to\R$ is the terminal/initial cost.

%% ==========================================================================
\section{The Cole-Hopf Transformation}
%% ==========================================================================

\subsection{\texorpdfstring{Exact case: $H(\bp) = \tfrac{1}{2}|\bp|^2$}{Exact case: H(p) = (1/2)|p|²}}

Define $\omega^\delta = e^{-\bv^\delta/\delta}$.  Then:
\begin{align}
  \omega_t^\delta &= -\frac{1}{\delta}\,\omega^\delta\, \bv_t^\delta, \\
  D\omega^\delta &= -\frac{1}{\delta}\,\omega^\delta\, D\bv^\delta, \\
  \Delta\omega^\delta &= \frac{1}{\delta^2}\,\omega^\delta\,|D\bv^\delta|^2
    - \frac{1}{\delta}\,\omega^\delta\,\Delta\bv^\delta.
\end{align}

From~\eqref{eq:hjvisc} with $H = \tfrac{1}{2}|D\bv^\delta|^2$:
\[
  \bv_t^\delta = \frac{\delta}{2}\Delta\bv^\delta - \frac{1}{2}|D\bv^\delta|^2.
\]

Substituting:
\begin{align}
  \omega_t^\delta
  &= -\frac{1}{\delta}\,\omega^\delta
     \Bigl(\frac{\delta}{2}\Delta\bv^\delta - \frac{1}{2}|D\bv^\delta|^2\Bigr)
  = -\frac{1}{2}\omega^\delta\Delta\bv^\delta
    + \frac{1}{2\delta}\omega^\delta|D\bv^\delta|^2, \\
  \frac{\delta}{2}\Delta\omega^\delta
  &= \frac{1}{2\delta}\omega^\delta|D\bv^\delta|^2
     - \frac{1}{2}\omega^\delta\Delta\bv^\delta.
\end{align}

Therefore $\omega_t^\delta = \frac{\delta}{2}\Delta\omega^\delta$, i.e.,
$\omega^\delta$ satisfies the \textbf{heat equation exactly}.

\begin{proposition}[Exact Cole-Hopf]\label{prop:exact}
  If $H(\bp) = \frac{1}{2}|\bp|^2$, then $\omega^\delta = e^{-\bv^\delta/\delta}$
  satisfies
  \[
    \omega_t^\delta = \frac{\delta}{2}\,\Delta\omega^\delta
    \quad\text{in } \R^n\times(0,T],
    \qquad
    \omega^\delta(0;\bx) = e^{-g(\bx)/\delta}.
  \]
\end{proposition}

\subsection{General case: the residual}

For a general Hamiltonian, define $c(t,\bx) := \frac{2\,H(t;\bx,D\bv^\delta)}{\delta\,|D\bv^\delta|^2}$
and set $\omega^\delta = e^{-c\,\bv^\delta}$.  Writing $f = -c\,\bv^\delta$,
the chain rule gives:
\begin{align}
  \omega_t^\delta &= \omega^\delta(-c_t\,\bv^\delta - c\,\bv_t^\delta), \label{eq:omega_t_general}\\
  \Delta\omega^\delta &= \omega^\delta\bigl(|Df|^2 + \Delta f\bigr), \label{eq:laplacian_omega}
\end{align}
where $Df = -Dc\,\bv^\delta - c\,D\bv^\delta$ and
$\Delta f = -\Delta c\,\bv^\delta - 2\langle Dc, D\bv^\delta\rangle - c\,\Delta\bv^\delta$.

\begin{proof}[Derivation of~\eqref{eq:laplacian_omega}]
We compute $\Delta\omega^\delta$ for $\omega^\delta = e^f$ where $f(\bx) = -c\,\bv^\delta$.

\medskip
\noindent\textbf{Step 1 (Gradient of $\omega^\delta$).}
By the chain rule applied to $\omega^\delta = e^f$:
\[
  \frac{\partial\omega^\delta}{\partial x_i}
  = e^f\,\frac{\partial f}{\partial x_i},
  \qquad\text{i.e.,}\qquad
  D\omega^\delta = \omega^\delta\,Df.
\]

\noindent\textbf{Step 2 (Second partial derivatives).}
Differentiate $\partial_{x_i}\omega^\delta = e^f\,\partial_{x_i}f$ once more using the product rule:
\[
  \frac{\partial^2\omega^\delta}{\partial x_i^2}
  = \frac{\partial}{\partial x_i}\!\Bigl(e^f\,\frac{\partial f}{\partial x_i}\Bigr)
  = \underbrace{e^f\,\Bigl(\frac{\partial f}{\partial x_i}\Bigr)^{\!2}}_{\text{chain rule on }e^f}
    \;+\;\underbrace{e^f\,\frac{\partial^2 f}{\partial x_i^2}}_{\text{product rule remainder}}.
\]

\noindent\textbf{Step 3 (Sum over all dimensions).}
The Laplacian $\Delta\omega^\delta = \sum_{i=1}^{n}\partial^2_{x_i}\omega^\delta$ gives:
\[
  \Delta\omega^\delta
  = \sum_{i=1}^{n}e^f\!\Bigl[\Bigl(\frac{\partial f}{\partial x_i}\Bigr)^{\!2}
    + \frac{\partial^2 f}{\partial x_i^2}\Bigr]
  = e^f\!\Biggl[\,\underbrace{\sum_{i=1}^n\Bigl(\frac{\partial f}{\partial x_i}\Bigr)^{\!2}}_{=\,|Df|^2}
    \;+\;\underbrace{\sum_{i=1}^n\frac{\partial^2 f}{\partial x_i^2}}_{=\,\Delta f}\Biggr]
  = \omega^\delta\bigl(|Df|^2 + \Delta f\bigr).
\]

\noindent\textbf{Step 4 (Expand $Df$ and $\Delta f$ via the product rule).}
Since $f = -c\,\bv^\delta$ and $c = c(t,\bx)$ is itself spatially varying:
\begin{align*}
  Df &= -Dc\,\bv^\delta - c\,D\bv^\delta, \\
  \Delta f &= -\Delta c\,\bv^\delta
    - 2\langle Dc, D\bv^\delta\rangle
    - c\,\Delta\bv^\delta,
\end{align*}
where the cross-term $-2\langle Dc, D\bv^\delta\rangle$ arises from the
product-rule identity $\Delta(c\,\bv^\delta) = (\Delta c)\,\bv^\delta
+ 2\langle Dc, D\bv^\delta\rangle + c\,\Delta\bv^\delta$.
\end{proof}

Substituting these expressions and using the HJ PDE
$\bv_t^\delta = \frac{\delta}{2}\Delta\bv^\delta - H$ yields the following
residual equation.

\begin{theorem}[Residual equation]\label{thm:residual}
  The transformed variable $\omega^\delta = e^{-c\,\bv^\delta}$ satisfies
  \begin{equation}\label{eq:residual}
    \omega_t^\delta - \frac{\delta}{2}\,\Delta\omega^\delta
    = \omega^\delta\cdot R(t,\bx),
  \end{equation}
  where the \textbf{residual} is
  \begin{equation}\label{eq:R}
    R = -c_t\,\bv^\delta
      + \delta\langle Dc, D\bv^\delta\rangle(1 - c\,\bv^\delta)
      - \frac{\delta}{2}\,(\bv^\delta)^2\,|Dc|^2
      + \frac{\delta}{2}\,\bv^\delta\,\Delta c.
  \end{equation}
  The leading-order terms ($cH - \frac{\delta}{2}c^2|D\bv^\delta|^2$) cancel
  by the choice of $c$; only terms involving derivatives of $c$ remain.
\end{theorem}

\begin{proof}
We set $f = -c\,\bv^\delta$ so that $\omega^\delta = e^f$, and compute
$\omega_t^\delta - \frac{\delta}{2}\Delta\omega^\delta$ in full.  The
derivation proceeds in six steps.

\medskip
\noindent\textbf{Step 1 (Factor out $\omega^\delta$).}
From~\eqref{eq:omega_t_general} and~\eqref{eq:laplacian_omega}:
\[
  \omega_t^\delta - \frac{\delta}{2}\,\Delta\omega^\delta
  = \omega^\delta\!\Bigl[
      f_t
      - \frac{\delta}{2}\bigl(|Df|^2 + \Delta f\bigr)
    \Bigr],
\]
so the residual is $R = f_t - \frac{\delta}{2}\,|Df|^2 - \frac{\delta}{2}\,\Delta f$.
We now expand each of these three pieces.

\medskip
\noindent\textbf{Step 2 (Compute $f_t$).}
Since $f = -c\,\bv^\delta$ with $c = c(t,\bx)$:
\[
  f_t = -c_t\,\bv^\delta - c\,\bv_t^\delta.
\]
Substitute the HJ PDE $\bv_t^\delta = \frac{\delta}{2}\,\Delta\bv^\delta - H$:
\begin{equation}\label{eq:ft_expanded}
  f_t
  = -c_t\,\bv^\delta - c\Bigl(\frac{\delta}{2}\,\Delta\bv^\delta - H\Bigr)
  = -c_t\,\bv^\delta
    - \frac{c\delta}{2}\,\Delta\bv^\delta
    + c\,H.
\end{equation}

\medskip
\noindent\textbf{Step 3 (Compute $|Df|^2$).}
From $Df = -Dc\,\bv^\delta - c\,D\bv^\delta$, expand the square:
\begin{equation}\label{eq:Df_sq}
  |Df|^2
  = |Dc|^2(\bv^\delta)^2
    + 2\,c\,\bv^\delta\,\langle Dc, D\bv^\delta\rangle
    + c^2\,|D\bv^\delta|^2.
\end{equation}

\medskip
\noindent\textbf{Step 4 (Compute $\Delta f$).}
From the product-rule identity
$\Delta(c\,\bv^\delta) = (\Delta c)\,\bv^\delta
  + 2\langle Dc, D\bv^\delta\rangle + c\,\Delta\bv^\delta$:
\begin{equation}\label{eq:Delta_f}
  \Delta f = -\Delta(c\,\bv^\delta)
  = -\Delta c\,\bv^\delta
    - 2\langle Dc, D\bv^\delta\rangle
    - c\,\Delta\bv^\delta.
\end{equation}

\medskip
\noindent\textbf{Step 5 (Assemble all terms).}
Substitute \eqref{eq:ft_expanded}--\eqref{eq:Delta_f} into
$R = f_t - \frac{\delta}{2}|Df|^2 - \frac{\delta}{2}\Delta f$:
\begin{align*}
  R &= \underbrace{
    \Bigl(-c_t\,\bv^\delta
    - \frac{c\delta}{2}\,\Delta\bv^\delta
    + c\,H\Bigr)}_{f_t} \\
  &\quad - \frac{\delta}{2}\,
    \underbrace{\Bigl(
      |Dc|^2(\bv^\delta)^2
      + 2\,c\,\bv^\delta\langle Dc, D\bv^\delta\rangle
      + c^2|D\bv^\delta|^2
    \Bigr)}_{|Df|^2} \\
  &\quad - \frac{\delta}{2}\,
    \underbrace{\Bigl(
      -\Delta c\,\bv^\delta
      - 2\langle Dc, D\bv^\delta\rangle
      - c\,\Delta\bv^\delta
    \Bigr)}_{\Delta f}.
\end{align*}

Distribute and list all seven resulting terms:
\begin{alignat*}{2}
  &\text{(a)}\;\; &&{-c_t\,\bv^\delta}, \\
  &\text{(b)}\;\; &&{-\tfrac{c\delta}{2}\,\Delta\bv^\delta}, \\
  &\text{(c)}\;\; &&{+c\,H}, \\
  &\text{(d)}\;\; &&{-\tfrac{\delta}{2}\,c^2|D\bv^\delta|^2}, \\
  &\text{(e)}\;\; &&{-c\delta\,\bv^\delta\langle Dc, D\bv^\delta\rangle}, \\
  &\text{(f)}\;\; &&{-\tfrac{\delta}{2}\,|Dc|^2(\bv^\delta)^2}, \\
  &\text{(g)}\;\; &&{+\delta\langle Dc, D\bv^\delta\rangle}, \\
  &\text{(h)}\;\; &&{+\tfrac{\delta}{2}\,\Delta c\,\bv^\delta}, \\
  &\text{(i)}\;\; &&{+\tfrac{c\delta}{2}\,\Delta\bv^\delta}.
\end{alignat*}

\medskip
\noindent\textbf{Step 6 (Cancellations and simplification).}

\emph{(I) The $\Delta\bv^\delta$ terms cancel.}
Terms (b) and (i): $\;-\frac{c\delta}{2}\Delta\bv^\delta + \frac{c\delta}{2}\Delta\bv^\delta = 0$.

\emph{(II) The leading-order terms cancel by the choice of $c$.}
From $c = \frac{2H}{\delta|D\bv^\delta|^2}$, compute:
\[
  \frac{\delta}{2}\,c^2|D\bv^\delta|^2
  = \frac{\delta}{2}\cdot\frac{4H^2}{\delta^2|D\bv^\delta|^4}\cdot|D\bv^\delta|^2
  = \frac{2H^2}{\delta|D\bv^\delta|^2}
  = c\,H.
\]
Therefore terms (c) and (d): $\;c\,H - \frac{\delta}{2}\,c^2|D\bv^\delta|^2 = 0$.

\emph{(III) Factor the $\langle Dc, D\bv^\delta\rangle$ terms.}
Terms (e) and (g):
\[
  \delta\langle Dc, D\bv^\delta\rangle
  - c\delta\,\bv^\delta\langle Dc, D\bv^\delta\rangle
  = \delta\langle Dc, D\bv^\delta\rangle\bigl(1 - c\,\bv^\delta\bigr).
\]

\emph{(IV) Collect the surviving terms:} (a), (III), (f), (h):
\[
  \boxed{
    R = -c_t\,\bv^\delta
      + \delta\langle Dc, D\bv^\delta\rangle(1 - c\,\bv^\delta)
      - \frac{\delta}{2}\,(\bv^\delta)^2|Dc|^2
      + \frac{\delta}{2}\,\bv^\delta\,\Delta c.
  }
\]
Every surviving term involves a derivative of $c$ ($c_t$, $Dc$, or $\Delta c$).
When $c$ is constant these all vanish and $R\equiv 0$, recovering the exact
Cole-Hopf result.
\end{proof}

\begin{corollary}
  $R \equiv 0$ if and only if $c$ is constant, which holds when
  $H/|D\bv^\delta|^2$ is independent of $(\bx,t)$.  For $H=\frac{1}{2}|\bp|^2$,
  $c = 1/\delta = \mathrm{const}$, recovering Proposition~\ref{prop:exact}.
\end{corollary}

\begin{remark}[Errata in the ICML submission]
  (i) Proposition~3.1 in the submission states the linearized equation as
  $\partial_t\omega^\delta - \frac{\delta}{2}\Delta\omega^\delta = H$;
  the RHS should be $0$ (for the exact case) or $\omega^\delta R$
  (for the general case).
  (ii) Equations~(11) and~(12) use $|Dg|^2/H^\delta$ as a prefactor;
  this equals $|D\bv^\delta|^2/H$ only at $t=0$.  For $t>0$, the
  correct prefactor involves $|D\bv^\delta|^2$ and $H(t;\bx,D\bv^\delta)$
  at the current time.
  (iii) Lemmas~3.3 and~3.4 are formally correct as identities but are
  not self-contained: they involve $D^2\bv^\delta$, $D_t D\bv^\delta$,
  $DH$, and $D_t H$, which are themselves unknown.
\end{remark}

%% ==========================================================================
\section{Quasi-Linearization Algorithm}
%% ==========================================================================

For general $H$, we propose an iterative scheme:

\medskip
\noindent\textbf{Algorithm} (Quasi-Linearized Cole-Hopf).

\begin{enumerate}
  \item \textbf{Initialize:} $\bv^{(0)}(t;\bx) = g(\bx)$, \;
        $c^{(0)} = \frac{2\,H(t;\bx,Dg)}{{\delta\,|Dg|^2}}$.
  \item \textbf{For $k = 0,1,2,\ldots$:}
    \begin{enumerate}
      \item Freeze $c^{(k)}$ at the current iterate.
      \item Solve the heat equation
            $\omega_t = \frac{\delta}{2}\Delta\omega$
            with initial data
            $\omega^{(k)}(0;\bx) = e^{-c^{(k)} g(\bx)}$.
      \item Recover $\bv^{(k+1)} = -(1/c^{(k)})\log\omega^{(k+1)}$.
      \item Update $D\bv^{(k+1)}$ and $c^{(k+1)} = \frac{2H(t;\bx,D\bv^{(k+1)})}{\delta|D\bv^{(k+1)}|^2}$.
      \item Check convergence: $\|\bv^{(k+1)} - \bv^{(k)}\|/\|\bv^{(k)}\| < \varepsilon$.
    \end{enumerate}
\end{enumerate}

%% ==========================================================================
\section{Gaussian Expectation Formulas}
%% ==========================================================================

At each iteration, the heat equation is solved via the Green's function:

\begin{lemma}[Value function recovery]\label{lem:value}
  Given the heat equation with initial data $\omega(0;\bx) = e^{-c\,g(\bx)}$,
  the value function at time $t$ is
  \begin{equation}\label{eq:value}
    \bv^\delta(t;\bx)
    = -\frac{1}{c}\,\log\E_{\bz\sim\mathcal{N}(0,I_n)}
      \bigl[e^{-c\,g(\bx + \sigma\bz)}\bigr],
    \qquad \sigma = \sqrt{\delta(T-t)}.
  \end{equation}
  For numerical stability, use the log-sum-exp identity:
  $-\frac{1}{c}\log\frac{1}{J}\sum_{j=1}^J e^{-c\,g(\bx+\sigma\bz_j)}$
  with $\bz_j\stackrel{\text{iid}}{\sim}\mathcal{N}(0,I_n)$.
\end{lemma}

\begin{lemma}[Gradient recovery]\label{lem:gradient}
  The spatial gradient is recovered as an importance-weighted ratio:
  \begin{equation}\label{eq:gradient}
    D\bv^\delta(t;\bx)
    = \frac{\E_{\bz}\bigl[Dg(\bx+\sigma\bz)\,e^{-c\,g(\bx+\sigma\bz)}\bigr]}
           {\E_{\bz}\bigl[e^{-c\,g(\bx+\sigma\bz)}\bigr]}.
  \end{equation}
\end{lemma}

\begin{proof}
  From $\omega = e^{-c\bv}$ and the heat-kernel representation
  $\omega(t;\bx) = \E[e^{-c\,g(\bx+\sigma\bz)}]$:
  \[
    D\omega = \E\bigl[-c\,Dg(\by)\,e^{-c\,g(\by)}\bigr]
    \cdot\frac{1}{\sigma}\,\frac{\partial\by}{\partial\bx}
    = \E\bigl[Dg(\by)\cdot(-c)\,e^{-cg(\by)}\bigr].
  \]
  Wait --- more carefully:
  $D_\bx\omega(t;\bx) = \E[D_\bx e^{-cg(\bx+\sigma\bz)}] = -c\,\E[Dg(\by)\,e^{-cg(\by)}]$.
  Since $\omega = e^{-c\bv}$ and $D\omega = -c\,e^{-c\bv}\,D\bv = -c\,\omega\,D\bv$,
  we get $D\bv = D\omega / (-c\,\omega)$, which gives~\eqref{eq:gradient}.
\end{proof}

%% ==========================================================================
\section{Error Analysis}
%% ==========================================================================

\subsection{Quasi-linearization error}

The residual $R$ in~\eqref{eq:R} involves $Dc$, $\Delta c$, $c_t$.
Since $c = 2H/(\delta|D\bv|^2)$, we have $Dc = O(D^2\bv, DH)$.
For smooth solutions with bounded Hessian $\|D^2\bv\| \le M_2$:
\[
  \|R\|_\infty \le C(\delta, M_2, \|DH\|, \|\bv\|_\infty)\cdot\|Dc\|_\infty.
\]

When the iteration converges ($\bv^{(k+1)}\approx\bv^{(k)}$), the
frozen coefficient $c^{(k)}$ stabilizes and $R\to 0$ effectively.

\subsection{Monte Carlo error}

The standard MC estimator for~\eqref{eq:value} converges at rate $O(1/\sqrt{J})$
where $J$ is the number of samples.  The effective variance depends on
$c$ and the range of $g$: when $c\,\|g\|_\infty \gg 1$ (small $\delta$),
the exponential weights become concentrated and the effective sample size
shrinks.  Importance sampling with a tilted proposal (e.g., Laplace
approximation around the mode of $e^{-cg}$) can reduce variance.

\subsection{Viscosity approximation error}

By Crandall-Lions (1984), $\|\bv^\delta - \bv\|_\infty \le k\sqrt{\delta}$
where $\bv$ is the inviscid viscosity solution.  Thus smaller $\delta$
improves the approximation quality but increases the MC variance
(the weights become more peaked).  This creates a fundamental
bias-variance tradeoff controlled by $\delta$.

\end{document}
